\documentclass[11pt, a4paper]{article}
\usepackage[a4paper, top=2.5cm, bottom=2.5cm, left=2cm, right=2cm]{geometry}
\usepackage{fontspec}
\usepackage[english, bidi=basic, provide=*]{babel}
\babelprovide[import, onchar=ids fonts]{english}
%\babelfont{rm}{Noto Sans}
\babelfont{rm}{TeX Gyre Termes}
\babelfont{sf}{TeX Gyre Heros}
\babelfont{tt}{TeX Gyre Cursor}

\usepackage{enumitem}
\usepackage{amsmath}

\title{\textbf{Chapter 1: Evolution, the Themes of Biology, and Scientific Inquiry \\ Exam Notes}}
\author{}
\date{}

\begin{document}

\maketitle

\section*{1.1 The Study of Life Reveals Unifying Themes}
Biology is the scientific study of life. The vast scope of biology is organized into five unifying themes.

\subsection*{Theme 1: Organization}
\begin{itemize}
    \item \textbf{Hierarchy of Life:} Unfolds from the global to the microscopic scale: Biosphere $\rightarrow$ Ecosystem $\rightarrow$ Community $\rightarrow$ Population $\rightarrow$ Organism $\rightarrow$ Organ $\rightarrow$ Tissue $\rightarrow$ Cell $\rightarrow$ Organelle $\rightarrow$ Molecule.
    \item \textbf{Emergent Properties:} Novel properties that arise at each upward level of the biological hierarchy due to the arrangement and interactions of parts as complexity increases.
    \item \textbf{Reductionism vs. Systems Biology:} 
    \begin{itemize}
        \item \textit{Reductionism} reduces complex systems to simpler components for study (e.g., studying molecular structure of DNA). 
        \item \textit{Systems Biology} models the dynamic behavior of whole biological systems based on interactions among parts.
    \end{itemize}
    \item \textbf{Structure and Function:} At every level, form correlates with function (e.g., a bird's wing structure allows flight).
    \item \textbf{The Cell:} The lowest level of organization capable of performing all activities required for life.
    \begin{itemize}
        \item \textit{Prokaryotic Cells:} Lack a nucleus and membrane-enclosed organelles (e.g., Bacteria, Archaea). Generally smaller.
        \item \textit{Eukaryotic Cells:} Contain a nucleus and membrane-bound organelles (e.g., Plants, Animals, Fungi).
    \end{itemize}
\end{itemize}

\subsection*{Theme 2: Information}
\begin{itemize}
    \item \textbf{DNA (Deoxyribonucleic Acid):} The genetic material contained within chromosomes. Inherited from parents, it directs development.
    \item \textbf{Gene Expression:} The process by which information in a gene directs the synthesis of a cellular product. DNA $\rightarrow$ transcribed into mRNA $\rightarrow$ translated into a linked series of amino acids $\rightarrow$ folds into a specific protein.
    \item \textbf{Universality of Genetic Code:} All forms of life use essentially the same genetic code, implying common ancestry.
    \item \textbf{Genomics and Proteomics:} 
    \begin{itemize}
        \item \textit{Genomics:} The study of whole sets of genes and their interactions within a species, and genome comparisons between species.
        \item \textit{Proteomics:} The study of sets of proteins and their properties.
        \item \textit{Bioinformatics:} The use of computational tools to process large volumes of genomic/proteomic data.
    \end{itemize}
\end{itemize}

\subsection*{Theme 3: Energy and Matter}
\begin{itemize}
    \item \textbf{Energy Flow:} Energy flows \textbf{one-way} through an ecosystem, usually entering as light (sun) and exiting as heat. Producers (plants) convert light into chemical energy; consumers feed on producers.
    \item \textbf{Chemical Cycling:} Matter is \textbf{recycled} within an ecosystem. Chemicals pass from plants to consumers to decomposers, and back to the environment.
\end{itemize}

\subsection*{Theme 4: Interactions}
\begin{itemize}
    \item \textbf{Feedback Regulation:} The mechanism by which biological processes self-regulate.
    \begin{itemize}
        \item \textit{Negative Feedback:} The most common form; the response reduces the initial stimulus (e.g., high blood glucose triggers insulin, which lowers blood glucose).
        \item \textit{Positive Feedback:} The end product speeds up its own production (e.g., blood clotting).
    \end{itemize}
    \item \textbf{Ecosystem Interactions:} Organisms interact continuously with their environments and other organisms. Human impacts, such as $CO_2$ emissions, lead to severe consequences like global climate change.
\end{itemize}

\section*{1.2 The Core Theme: Evolution}
Evolution accounts for the unity (shared features) and diversity (differences) of life on Earth.

\subsection*{Classifying the Diversity of Life}
\begin{itemize}
    \item \textbf{Three Domains of Life:}
    \begin{enumerate}
        \item \textbf{Bacteria:} Most diverse/widespread prokaryotes.
        \item \textbf{Archaea:} Prokaryotes, often living in extreme environments.
        \item \textbf{Eukarya:} All eukaryotes. Includes three multicellular kingdoms defined by nutrition:
        \begin{itemize}
            \item \textit{Plantae:} Produce their own food (photosynthesis).
            \item \textit{Fungi:} Absorb nutrients from surroundings.
            \item \textit{Animalia:} Ingest other organisms.
            \item \textit{Protists:} Mostly unicellular, highly diverse eukaryotes.
        \end{itemize}
    \end{enumerate}
\end{itemize}

\subsection*{Charles Darwin and Natural Selection}
\begin{itemize}
    \item \textbf{Descent with Modification:} Darwin's phrase for evolution. Current species arise from a succession of ancestors that differed from them.
    \item \textbf{Natural Selection:} The mechanism for evolutionary adaptation.
    \begin{itemize}
        \item \textit{Observations:} (1) Individuals vary in heritable traits; (2) Overproduction of offspring leads to competition; (3) Species suit their environments.
        \item \textit{Inference:} Individuals with advantageous traits are more likely to survive and reproduce. Over time, favorable traits accumulate in a population.
    \end{itemize}
    \item \textbf{Tree of Life:} Represents evolutionary relationships. Adaptive radiation (like Galápagos finches branching from a common ancestor) highlights how populations adapt to different environments over time.
\end{itemize}

\section*{1.3 Scientific Inquiry}
Science is an approach to understanding the natural world based on inquiry.

\subsection*{Reasoning and Data}
\begin{itemize}
    \item \textbf{Data:} Recorded observations. \textit{Qualitative} (descriptions) vs. \textit{Quantitative} (numerical measurements).
    \item \textbf{Inductive Reasoning:} Deriving a general conclusion from specific, repeated observations (specific $\rightarrow$ general).
    \item \textbf{Deductive Reasoning:} Using general premises to extrapolate specific predictions (general $\rightarrow$ specific). Uses "If... then" logic.
\end{itemize}

\subsection*{Hypotheses and Experiments}
\begin{itemize}
    \item \textbf{Hypothesis:} An explanation based on observations/assumptions that leads to a testable prediction. Must be testable and falsifiable. Science cannot test supernatural explanations.
    \item \textbf{Controlled Experiment:} Compares an experimental group to a control group, usually varying only one factor.
    \begin{itemize}
        \item \textit{Independent Variable:} The factor manipulated by the researcher.
        \item \textit{Dependent Variable:} The factor measured to determine the effect of the manipulation.
    \end{itemize}
\end{itemize}

\subsection*{Scientific Theories}
\begin{itemize}
    \item A scientific theory is much broader in scope than a hypothesis, generates new hypotheses, and is supported by a massive body of evidence (e.g., Theory of Natural Selection).
\end{itemize}

\section*{1.4 Cooperative Approach and Science vs. Society}
\begin{itemize}
    \item \textbf{Science vs. Technology:} Science aims to understand natural phenomena; technology applies scientific knowledge for a specific purpose.
    \item \textbf{Model Organisms:} Species chosen for research (like \textit{E. coli}, mice, or fruit flies) because they are easy to grow in labs and offer generalizable insights.
    \item \textbf{Diversity:} Science benefits immensely from the integration of diverse viewpoints, backgrounds, and collaborative peer review.
\end{itemize}

\end{document}
